\documentclass[11pt]{article}
\usepackage[T1]{fontenc}
\usepackage[margin=1in]{geometry}
\usepackage{amsmath,amssymb,amsthm}
\usepackage{hyperref}
\setlength{\emergencystretch}{3em}
\newtheorem{theorem}{Theorem}
\newtheorem{lemma}[theorem]{Lemma}
\newtheorem{proposition}[theorem]{Proposition}
\newcommand{\R}{\mathbb R}
\newcommand{\abs}[1]{\lvert#1\rvert}
\newcommand{\norm}[1]{\lVert#1\rVert_1}
\title{A sharp three-halves Erd\H{o}s--Mordell bound\\for arbitrary taxicab triangles}
\author{Myeongjin Kang}
\date{26 October 2026 --- research draft}
\begin{document}
\maketitle
\begin{abstract}
For a point in an arbitrary triangle in the taxicab plane, we prove that the sum of the distances to the vertices is at least three halves the sum of the distances to the supporting lines of its sides. No restriction on the orientation of the triangle is imposed. A barycentric comparison reduces the problem to a scalar inequality, which follows by concavity on a three-vertex simplex. The constant is sharp. The result addresses Conjecture~1 of Petrovi\'c, Male\v{s}evi\'c and Banjac, recorded as MathDB~341671. The proof is accompanied by Lean~4.34.1 sources. Prior-art and final certificate review remain in progress; this draft is not yet counted as an accepted new discovery.
\end{abstract}
\noindent\textit{Provenance.} AI-assisted research. The source conjecture and its previously proved restricted-orientation case are due to Petrovi\'c, Male\v{s}evi\'c and Banjac. The scalar argument and its integration are developed in this research packet. Original prose is CC~BY~4.0 and original code is MIT.

\section{Exact geometric statement}
For $X=(x_1,x_2)\in\R^2$, write $\norm X=\abs{x_1}+\abs{x_2}$. For a line $L$ with equation $a x+b y+c=0$, where $(a,b)\ne(0,0)$, use the taxicab distance
\begin{equation}\label{eq:line}
 d_1(P,L)=\frac{\abs{aP_1+bP_2+c}}{\max(\abs a,\abs b)}.
\end{equation}
This is the supporting-line convention explicitly used in equations~(5)--(6) of the primary paper~\cite{pmb}. In particular, no Euclidean perpendicular distance is substituted.

Let $A,B,C$ be distinct vertices, let $P$ belong to their closed convex hull, and put
\[
 \begin{gathered}
 R_A=\norm{A-P},\quad R_B=\norm{B-P},\quad R_C=\norm{C-P},\\
 r_a=d_1(P,BC),\quad r_b=d_1(P,CA),\quad r_c=d_1(P,AB).
 \end{gathered}
\]
Here $BC$ means the full supporting line, and similarly for the other sides.
\begin{theorem}\label{thm:main}
For every such configuration,
\begin{equation}\label{eq:main}
 R_A+R_B+R_C\ \ge\ \frac32(r_a+r_b+r_c).
\end{equation}
Consequently this holds for every interior point of every nondegenerate taxicab triangle. The constant $3/2$ cannot be increased, even if only nondegenerate triangles and strictly interior points are considered.
\end{theorem}
The proof is invariant under translations and simultaneous relabelling of vertices and opposite sides. It never uses rotation as a taxicab isometry.

\section{A scalar inequality}
\begin{lemma}\label{lem:scalar}
Suppose $A,B,C>0$ and $u,v,w\ge0$ satisfy
\[
 u\le\min(B,C),\qquad v\le\min(A,C),\qquad w\le\min(A,B),
\]
and
\begin{equation}\label{eq:F}
 F(u,v,w):=\frac{u}{A+u}+\frac{v}{B+v}+\frac{w}{C+w}\le1.
\end{equation}
Then $3(u+v+w)\le2(A+B+C)$.
\end{lemma}
\begin{proof}
Relabel so that $A\le B\le C$, relabelling $u,v,w$ with their corresponding denominators. Their upper bounds reduce to $(B,A,A)$. Put $t=2(A+B+C)/3$. If $u+v+w>t$, multiply $u,v,w$ by $t/(u+v+w)\in(0,1)$. The scaled point stays in the box and lies on the slice with sum $t$. Each function $x/(a+x)$, $a>0$, is strictly increasing on $[0,\infty)$; at least one coordinate was positive. Thus the scaled value of $F$ is strictly smaller than its original value, and in particular is strictly less than one.

We show instead that $F\ge1$ throughout this slice. Set $D=B+2A-t$. Feasibility gives $D\ge0$. Since
\[
 t\ge A+B,\qquad t\ge2A,
\]
we have $D\le A\le B$. If $D=0$, the slice is the point $(B,A,A)$, whose first two summands in $F$ add to one. If $D>0$, the nonnegative deficits $(B-u,A-v,A-w)$ add to $D$. Dividing them by $D$ expresses any point of the slice as a convex combination of
\begin{align*}
 V_1&=(B-D,A,A)=(t-2A,A,A),\\
 V_2&=(B,A-D,A)=(B,t-A-B,A),\\
 V_3&=(B,A,A-D)=(B,A,t-A-B).
\end{align*}
Each displayed coordinate is nonnegative. Each function $x/(a+x)$ is concave there, so $F$ is concave. It is therefore enough to verify $F(V_i)\ge1$ at these three vertices.

At $V_1$, put $x=2(B+C-2A)/3$. Clearing positive denominators shows that
\[
 \frac{x}{A+x}+\frac{A}{B+A}+\frac{A}{C+A}\ge1
 \quad\Longleftrightarrow\quad x(2A+B+C)\ge BC-A^2.
\]
The difference between the last two sides is
\[
 \frac{2B^2+BC+2C^2-5A^2}{3}\ge0,
\]
since $B,C\ge A>0$.

At $V_2$, put $y=(2C-A-B)/3\ge0$. We have
\[
 \frac{y}{B+y}\ge\frac{C-B}{B+C}
 \ge\frac{A(C-B)}{(A+B)(A+C)}.
\]
Indeed the cross-product difference in the first comparison is
\[
 y(B+C)-(C-B)(B+y)=\frac{B(B+C-2A)}3\ge0.
\]
The second follows from $C\ge B$ and $(A+B)(A+C)-A(B+C)=A^2+BC>0$. The first and third summands of $F(V_2)$ add to
\[
 \frac{B}{A+B}+\frac{A}{C+A}
 =1-\frac{A(C-B)}{(A+B)(A+C)},
\]
proving $F(V_2)\ge1$.

At $V_3$, the first two summands add to $B/(A+B)+A/(B+A)=1$, and the third is nonnegative. This proves $F\ge1$ on the slice. It contradicts the strictly smaller value obtained by scaling a purported violating point, and proves the lemma.
\end{proof}

\begin{lemma}[Weighted form, including zero lengths]\label{lem:weighted}
Let $A,B,C,u,v,w\ge0$ satisfy the same upper bounds as Lemma~\ref{lem:scalar}. Suppose $a,b,c\ge0$, $a+b+c=1$, and
\[
 (1-a)u\le aA,\qquad(1-b)v\le bB,\qquad(1-c)w\le cC.
\]
Then $3(u+v+w)\le2(A+B+C)$.
\end{lemma}
\begin{proof}
For every $\varepsilon>0$, the weighted inequalities imply
\[
 \frac{u}{A+\varepsilon+u}\le a,\qquad
 \frac{v}{B+\varepsilon+v}\le b,\qquad
 \frac{w}{C+\varepsilon+w}\le c.
\]
Apply Lemma~\ref{lem:scalar} to $A+\varepsilon,B+\varepsilon,C+\varepsilon$ and the unchanged $u,v,w$. It gives
\[
 3(u+v+w)\le2(A+B+C)+6\varepsilon.
\]
If the desired conclusion failed, choosing $\varepsilon=[3(u+v+w)-2(A+B+C)]/12>0$ would give an immediate contradiction. Thus no analytic limiting theorem is required.
\end{proof}

\section{Barycentric comparison and the main proof}
\begin{lemma}\label{lem:geometry}
There exist nonnegative $a,b,c$ with sum one such that
\begin{align*}
 r_a&\le\min(R_B,R_C), &(1-a)r_a&\le aR_A,\\
 r_b&\le\min(R_A,R_C), &(1-b)r_b&\le bR_B,\\
 r_c&\le\min(R_A,R_B), &(1-c)r_c&\le cR_C.
\end{align*}
\end{lemma}
\begin{proof}
Choose barycentric coordinates $P=aA+bB+cC$ with $a,b,c\ge0$ and $a+b+c=1$, which exist precisely because $P$ is in the closed triangle. A direct coordinate inequality is
\[
 \abs{\alpha x+\beta y}\le\max(\abs\alpha,\abs\beta)(\abs x+\abs y).
\]
Thus every point $Q$ of the line $BC$ gives $r_a\le\norm{P-Q}$. Choosing the endpoints $B,C$ proves the two upper bounds on $r_a$.

For completeness, the weighted comparison can be proved without division by $1-a$, so it also covers a boundary point with $a=1$. Let $\ell(X)=\alpha X_1+\beta X_2+\gamma$ be the line function for $BC$. Its affine barycentric identity gives $\ell(P)=a\ell(A)$ because $\ell(B)=\ell(C)=0$. Therefore
\[
 (1-a)\ell(P)=a(\ell(A)-\ell(P)).
\]
Since $a,1-a\ge0$, taking absolute values, applying the coordinate inequality to $A-P$, and dividing by $\max(\abs\alpha,\abs\beta)>0$ gives $(1-a)r_a\le aR_A$. Cyclic permutation proves the other assertions.
\end{proof}
\begin{proof}[Proof of Theorem~\ref{thm:main}, inequality]
Apply Lemma~\ref{lem:weighted} with $(A,B,C)=(R_A,R_B,R_C)$ and $(u,v,w)=(r_a,r_b,r_c)$, using Lemma~\ref{lem:geometry}. This gives exactly~\eqref{eq:main}, including closed-triangle boundary points. Its restriction to the ordinary interior is the source conjecture.
\end{proof}

\section{Sharpness}
The upper bound $3/2$ for the optimal taxicab coefficient is prior art: Kurusa~\cite[Section~7.3.4]{kurusa} attributes a configuration of ratio $3/2$ to Ghandehari and Martini~\cite[Figure~3]{gm}. The family below supplies a transparent strictly interior approximation and is not claimed as a separate new result. The candidate contribution is the universal lower bound.

Take vertices $A=(0,1)$, $B=(-1,0)$, $C=(1,0)$ and strictly interior point $P=(0,\varepsilon)$, where $0<\varepsilon<1$. Their barycentric coordinates are $\varepsilon,(1-\varepsilon)/2,(1-\varepsilon)/2$, all positive. Directly from~\eqref{eq:line},
\[
 R_A+R_B+R_C=3+\varepsilon,\qquad
 r_a+r_b+r_c=2-\varepsilon.
\]
The ratio tends to $3/2$ as $\varepsilon\downarrow0$. More explicitly, for any $K>3/2$, choose
\[
 0<\varepsilon<\min\left(1,\frac{2K-3}{K+1}\right).
\]
Then $3+\varepsilon<K(2-\varepsilon)$, so $K$ is not a valid universal coefficient. This proves the sharpness assertion. The boundary point $\varepsilon=0$ attains equality in the closed-triangle version.

\section{Source fidelity, formal coverage, and novelty limits}
Petrovi\'c, Male\v{s}evi\'c and Banjac proved the coefficient $3/2$ for triangles in specified coordinate positions, and recorded the unrestricted assertion as Conjecture~1~\cite{pmb}. Their concluding conjecture remains in the published 2020 paper, page~1318. The present proof handles arbitrary coordinates directly, without using a rotation to place a side on a coordinate axis. It is a substantive geometric argument for the unrestricted inequality.

The certificate uses an explicit coordinate $L^1$ size and the exact supporting-line formula; it does not use Lean's default maximum norm on a product of two real lines. Its scalar core proves concavity, all three simplex-vertex inequalities, all six orderings of the vertex sizes, the scaling contradiction, and the weighted extension allowing zero lengths. The geometry module proves the line comparison and all cyclic barycentric bounds. The integration module targets the ordinary closed convex hull and its topological interior. Exact theorem names, dependency revisions, source hashes, and final audit state are supplied in the accompanying packet. All 55 local declarations, including the ordinary-interior theorem and sharpness, pass a fresh isolated replay against pinned Lean and mathlib dependencies. An independent replay also passed with only the three standard axioms. Independent mathematical and source reviews found no defect. Novelty remains on hold pending the closest prior-work comparison.

Prior-art investigations include the earlier taxicab paper of S\"onmez~\cite{sonmez}, the 2018 work of Ghandehari and Martini~\cite{gm}, and Kurusa's 2024 discussion of the Erd\H{o}s ratio~\cite{kurusa}. The full text of the 2018 paper has not yet been obtained from an accessible primary repository. Its relation to the present bound therefore remains an explicit review item. No claim of absolute priority is made, and no discovery count is assigned on the basis of unsuccessful web searches. The sharpness example is elementary and is not claimed as a separate new discovery. The proof also has an abstract normed-space interpretation, which is not included as a separately counted or fully formalized theorem here.

\clearpage
\begin{thebibliography}{9}
\bibitem{pmb} M. Petrovi\'c, B. Male\v{s}evi\'c and B. Banjac, \emph{On the Erd\H{o}s--Mordell inequality for triangles in taxicab geometry}, Journal of Mathematical Inequalities 14(4) (2020), 1299--1319. \href{https://doi.org/10.7153/jmi-2020-14-84}{doi:10.7153/jmi-2020-14-84}. Preprint \href{https://arxiv.org/abs/1901.08758}{arXiv:1901.08758}; Conjecture~1.
\bibitem{mathdb} MathDB, \emph{The Erd\H{o}s--Mordell inequality for arbitrary taxicab triangles}, entry~341671, accessed 26 October 2026. Exact API snapshot and hash accompany the certificate.
\bibitem{sonmez} N. S\"onmez, \emph{A taxicab version of the Erd\H{o}s--Mordell theorem}, Kragujevac Journal of Mathematics 32 (2009), 157--164. The later paper~\cite{pmb} discusses the failure of its coefficient-two claim under the supporting-line distance convention.
\bibitem{gm} M. Ghandehari and H. Martini, \emph{On the Erd\H{o}s--Mordell inequality for normed planes and spaces}, Studia Scientiarum Mathematicarum Hungarica 55(2) (2018), 174--189. \href{https://doi.org/10.1556/012.2018.55.2.1392}{doi:10.1556/012.2018.55.2.1392}.
\bibitem{kurusa} \'A. Kurusa, \emph{Metric Characterizations of Projective-Metric Spaces}, in A. Papadopoulos (ed.), Surveys in Geometry II (2024), 225--256, Section~7.3.4. \href{https://doi.org/10.1007/978-3-031-43510-2_7}{doi:10.1007/978-3-031-43510-2\_7}.
\end{thebibliography}
\end{document}
